% !PRINTED 510
\weberpage{510}

We can represent $\Phi(M,u)$ as an \gesperrt{integral function of} $u$, and can then replace this function by its greatest common divisor with $H_{-D}(u)$.

We may assume $\alpha$ and $\beta$ to be positive. For $\beta$ this follows from the assumptions already made. For $\alpha$ we can always achieve it by increasing $B$ by a multiple of $2A$, whereby we arrive at an equivalent (parallel) form. In this way $Bx+y$ can be made positive.

If we therefore determine the quotient
\[
\frac{\eta\left(\dfrac{\gamma+\delta\omega}{\alpha+\beta\omega}\right)}{\eta(\omega)} = E\binom{\alpha,\ \beta}{\gamma,\ \delta}
\]
by § 38, (15), then, since $\beta$ is even or odd according as $x$ is even or odd, it follows that:

\weq{12}{\begin{aligned}
x\equiv 0 \pmod 2:&\quad M=\left(\frac{\beta}{\alpha}\right)i^{3\frac{\alpha-1}{2}}e^{\frac{\pi i}{4}\alpha(\gamma-\beta)}\sqrt{\alpha+\beta\omega}^{\,3},\\
x\equiv 1 \pmod 2:&\quad M=\left(\frac{\alpha}{\beta}\right)i^{3\frac{1-\beta}{2}}e^{\frac{\pi i}{4}\beta(\alpha+\delta)}\sqrt{-i(\alpha+\beta\omega)}^{\,3}.
\end{aligned}}

The square roots
\weq{13}{\begin{aligned}
\sqrt{\alpha+\beta\omega}&=\sqrt{\dfrac{y+x\sqrt D}{2}},\\
\sqrt{-i(\alpha+\beta\omega)}&=\sqrt{-i\dfrac{y+x\sqrt D}{2}}
\end{aligned}}
have positive real part. These values of $M$ therefore satisfy equation (11). If equal ones occur among them, the degree of equation (11) can be lowered by rational means through the separation of multiple roots.

Since equation (11) is to be applied to the decomposition of the class equation according to the signs of $M$, it is of importance to decide whether among the values of $M$ in (12) belonging to the same $\omega$ there occur any that differ only in sign.

We investigate when two different ones among the values (12), say $M$, $M'$, have the same 8th power. If
\weq{14}{M^8=M'^8,}
then, if $\varrho$ is a twelfth root of unity and $x,y;\ x',y'$ are two different solutions of equation (8), we must have:
\weq{15}{y'+x'\sqrt D=\varrho(y+x\sqrt D).}
